\subsection{Contexto y problemática}

El crecimiento exponencial de los datos digitales y la adopción generalizada de 
aplicaciones distribuidas han puesto de manifiesto las limitaciones de los modelos 
tradicionales de almacenamiento centralizado. Las soluciones basadas en infraestructuras 
cloud, ampliamente utilizadas en la actualidad, presentan una fuerte dependencia de 
entidades proveedoras, lo que introduce problemas relacionados con la confianza, 
la disponibilidad de los datos y la posibilidad de manipulación o pérdida de 
información \cite{tanenbaum2007distributed}.

En este contexto, la integridad y la verificabilidad de los datos se convierten 
en requisitos fundamentales, especialmente en escenarios donde múltiples actores 
interactúan sin una relación de confianza previa. La necesidad de garantizar que 
un fichero no ha sido alterado, así como de poder demostrar su existencia en un 
momento determinado, resulta crítica en aplicaciones como la gestión documental, 
sistemas de auditoría o intercambio seguro de información \cite{nist2011cloud,iso27001}.

Sin embargo, los sistemas tradicionales no ofrecen mecanismos nativos 
que permitan asegurar estas propiedades de forma transparente y descentralizada. 
La verificación de la autenticidad de los datos suele depender de terceros 
de confianza o de infraestructuras centralizadas, lo que introduce puntos 
únicos de fallo y limita la trazabilidad de la 
información \cite{tanenbaum2007distributed}.

Ante esta situación, surge la necesidad de explorar nuevos paradigmas de 
almacenamiento y verificación que eliminen la dependencia de entidades centrales y 
permitan garantizar la integridad, disponibilidad y autenticidad de los datos de 
forma distribuida. Este problema constituye el punto de partida del presente trabajo.